\chapter{Réalisation et Validation Industrielle}
\section*{Introduction}
\addcontentsline{toc}{section}{Introduction}
Ce dernier chapitre décrit l'environnement de développement, l'architecture
de déploiement Docker, les tests de validation et les résultats de la
validation terrain sur 7 fermes pilotes tunisiennes.

% ── 4.1 Environnement de développement ──────────────────────
\section{Environnement de développement}

\begin{table}[H]
\centering
\caption{Environnement de développement}
\label{tab:env_dev}
\begin{tabularx}{\linewidth}{llX}
\toprule
\textbf{Niveau} & \textbf{Composant} & \textbf{Détails} \\
\midrule
Matériel & Station de dev & Intel Core i7-12700H, 32 GB RAM, GPU NVIDIA RTX 3060 \\
         & Appareils IoT  & 2$\times$ ESP32-WROOM-32, capteurs DHT22, HX711, YF-S201 \\
OS       & Développement  & Ubuntu 22.04 LTS + Windows 11 (WSL2) \\
         & Production     & Ubuntu 22.04 LTS (VPS) \\
Runtime  & Python         & 3.11.8 (venv isolé) \\
         & Node.js        & 20.11 LTS \\
         & Docker         & 25.0.3 + Compose v2.24 \\
IDE      & Backend        & VS Code + Pylance, Black, Ruff \\
         & Frontend       & VS Code + ESLint, Prettier \\
VCS      & Git            & GitHub (branche main + feature branches) \\
MLOps    & MLflow         & $\geq$2.10.0 (tracking server local) \\
         & DVC            & $\geq$3.45.0 (artefacts modèles) \\
\bottomrule
\end{tabularx}
\end{table}

% ── 4.2 Architecture de déploiement ─────────────────────────
\section{Architecture de déploiement Docker}

Le fichier \code{docker-compose.yml} orchestre 3 services actifs (+ Caddy optionnel HTTPS) :

\begin{figure}[H]
\centering
\begin{tikzpicture}[
  svc/.style={draw, rounded corners=4pt, minimum width=3.2cm, minimum height=0.8cm,
              align=center, font=\small\bfseries},
  port/.style={font=\tiny, gray},
  arrow/.style={->, thick, >=stealth},
  darrow/.style={<->, thick, >=stealth, gray!60},
]
% Internet / Client
\node[draw, rounded corners=8pt, minimum width=3.2cm, minimum height=0.8cm,
      fill=gray!10, font=\small\bfseries] (internet) at (0, 3.5) {Internet / Clients};

% Caddy (optional HTTPS)
\node[svc, fill=gray!15, draw=gray!50, dashed] (caddy) at (0, 2.1)
  {Caddy (opt.) \\ \textcolor{gray}{\tiny HTTPS :443}};

% Nginx frontend
\node[svc, fill=blue!15, draw=blue!50] (nginx) at (-3.5, 0.5)
  {nginx:alpine \\ \textcolor{blue!60}{\tiny :80 (React/Vite)}};

% FastAPI backend
\node[svc, fill=green!15, draw=green!50] (backend) at (0, 0.5)
  {python:3.11-slim \\ \textcolor{green!60!black}{\tiny :8000 (uvicorn ×2)}};

% PostgreSQL
\node[svc, fill=orange!15, draw=orange!50] (db) at (3.5, 0.5)
  {postgres:16-alpine \\ \textcolor{orange!70}{\tiny :5432}};

% ESP32 IoT
\node[draw, rounded corners, fill=teal!10, draw=teal!50, minimum width=2.8cm,
      minimum height=0.7cm, font=\scriptsize, align=center] (iot) at (3.5, -1.2)
  {ESP32 Node A+B \\ \texttt{\tiny smartfarm/\#} MQTT};

% Arrows
\draw[arrow] (internet) -- (caddy);
\draw[arrow] (caddy) -- (nginx);
\draw[arrow] (caddy) -- (backend);
\draw[darrow] (nginx) -- node[above, font=\tiny]{/api proxy} (backend);
\draw[arrow] (backend) -- node[above, font=\tiny]{SQLAlchemy} (db);
\draw[arrow] (iot) -- node[right, font=\tiny]{paho-mqtt} (backend);

% Volume annotation
\node[font=\tiny, gray] at (3.5, -0.3) {Volumes persistants};

\end{tikzpicture}
\caption{Architecture de déploiement Docker — 3 services actifs (Caddy optionnel), ports réels, flux de communication.
Backend : Uvicorn 2 workers, \code{ACCESS\_TOKEN\_EXPIRE\_MINUTES=10080} (7 jours).}
\label{fig:deployment_arch}
\end{figure}

Le fichier \code{docker-compose.yml} orchestre ces services :

\begin{table}[H]
\centering
\caption{Services Docker Compose en production (3 actifs + 1 optionnel)}
\label{tab:docker}
\begin{tabularx}{\linewidth}{llX}
\toprule
\textbf{Service} & \textbf{Image} & \textbf{Rôle} \\
\midrule
\code{db}          & postgres:16-alpine & BDD principale, volume persistant \code{postgres\_data} \\
\code{backend}     & Dockerfile (Python 3.11) & FastAPI + Uvicorn (2 workers), port 8000, MQTT HiveMQ \\
\code{frontend}    & Dockerfile (nginx:alpine) & React/Vite build, port 80, proxy \code{/api/v1 $\rightarrow$ backend} \\
\code{caddy}       & caddy:2-alpine     & HTTPS auto Let's Encrypt — \textit{(commenté par défaut)} \\
\bottomrule
\end{tabularx}
\par\smallskip
{\small \textit{Note :} Caddy activé si \code{DOMAIN=} configuré dans \code{.env.production}.}
\end{table}

Les dépendances de démarrage sont gérées par \code{healthcheck} :
\code{postgres} doit être \code{healthy} avant \code{backend} ;
\code{backend} doit être \code{healthy} avant \code{frontend}.

% ── 4.3 Interfaces et fonctionnalités ───────────────────────
\section{Interfaces réalisées}

\subsection{Dashboard principal (Owner)}

Le tableau de bord principal agrège :
\begin{itemize}
  \item KPIs temps réel : nombre d'animaux actifs, alertes non lues,
        température moyenne IoT, production du jour ;
  \item Graphiques de tendance (Recharts) : FCR sur 42 jours, courbe de
        mortalité cumulée, production d'œufs hebdomadaire ;
  \item Carte GIS (Leaflet/MapLibre) : localisation fermes, vétérinaires
        dans un rayon paramétrable ;
  \item Flux alertes WebSocket temps réel : badge animé, toasts critiques.
\end{itemize}

\subsection{Scanner IA (AIScanner)}

L'interface AIScanner permet :
\begin{itemize}
  \item Upload ou capture caméra directe ;
  \item Affichage des bounding boxes YOLO sur l'image ;
  \item Panel de réponse Darija avec niveau de confiance coloré ;
  \item Historique des diagnostics (DiagnosticHistory) scrollable.
\end{itemize}

\subsection{ERP Aviculture}

Le module ERP avicole comprend 6 vues fonctionnelles :
tableau de bord lot avec graphiques FCR/mortalité en temps réel,
formulaire d'enregistrement alimentation, journal de ponte,
registre sanitaire, ventes et inventaire stock.

\subsection{PWA Worker (mobile)}

L'application mobile Worker est installable sur Android/iOS via le
manifest PWA. Elle fonctionne \textit{intégralement hors ligne} grâce à :
\begin{itemize}
  \item Dexie IndexedDB pour le stockage local des tâches assignées ;
  \item Service Worker pour la mise en cache des assets statiques ;
  \item Sync différée : les actions terrain (rapport tâche, télémétrie manuelle)
        sont enqueued et synchronisées au rétablissement de la connexion.
\end{itemize}

% ── 4.4 Captures d'écran ─────────────────────────────────────
\section{Captures d'écran de l'application}

\subsection{Dashboard principal}

\begin{figure}[H]
\centering
\begin{tikzpicture}[font=\tiny]
  \draw[rounded corners=6pt,thick,blue!40,fill=gray!4](0,0)rectangle(13,8.5);
  \fill[blue!70](0,7.8)rectangle(13,8.5);
  \node[white,font=\small\bfseries]at(6.5,8.15){Smart Farm AI — Tableau de bord Owner};
  \foreach \i/\lbl/\val/\clr in {
    0/{Animaux actifs}/{342}/{green!60!black},
    1/{Alertes}/{3}/{red!70},
    2/{Temp.\ moy.}/{24.3\,°C}/{orange!80},
    3/{Production}/{48.5\,kg}/{teal!70}}{
    \colorlet{kpifillclr}{\clr!15}
    \colorlet{kpibordclr}{\clr}
    \fill[kpifillclr,draw=kpibordclr,rounded corners=3pt]({0.3+\i*3.18},6.8)rectangle({3.1+\i*3.18},7.6);
    \node[kpibordclr,font=\scriptsize\bfseries]at({1.7+\i*3.18},7.3){\val};
    \node at({1.7+\i*3.18},6.95){\lbl};}
  \fill[white,draw=gray!40,rounded corners=3pt](0.3,4.0)rectangle(6.3,6.6);
  \node[font=\scriptsize\bfseries]at(3.3,6.38){FCR — 42 jours};
  \draw[blue!60,thick](0.6,4.5)--(1.3,4.4)--(2.0,4.3)--(2.7,4.5)--(3.4,4.8)--(4.1,5.1)--(4.8,5.5)--(5.5,5.8)--(6.0,6.1);
  \fill[white,draw=gray!40,rounded corners=3pt](6.6,4.0)rectangle(12.7,6.6);
  \node[font=\scriptsize\bfseries]at(9.65,6.38){Télémétrie IoT temps réel};
  \draw[green!60!black,thick](6.9,4.5)--(7.5,4.7)--(8.1,4.6)--(8.7,4.4)--(9.3,4.8)--(9.9,5.0)--(10.5,4.8)--(11.1,4.7)--(11.7,4.9);
  \draw[orange!60,dashed](6.9,5.2)--(7.5,5.3)--(8.1,5.5)--(8.7,5.3)--(9.3,5.4)--(9.9,5.6)--(10.5,5.3)--(11.1,5.4)--(11.7,5.5);
  \fill[teal!8,draw=teal!40,rounded corners=3pt](0.3,0.5)rectangle(6.3,3.8);
  \node[font=\scriptsize\bfseries]at(3.3,3.6){Carte GIS — Fermes \& Vétérinaires};
  \fill[teal!30](1.6,2.0)circle(0.15);\fill[teal!30](2.9,2.4)circle(0.15);\fill[teal!30](3.8,1.7)circle(0.15);
  \fill[red!50](4.5,2.6)circle(0.12);\fill[red!50](2.1,1.3)circle(0.12);
  \fill[white,draw=gray!40,rounded corners=3pt](6.6,0.5)rectangle(12.7,3.8);
  \node[font=\scriptsize\bfseries]at(9.65,3.6){Alertes WebSocket};
  \foreach \y/\msg/\clr in {
    3.15/{[URGENT] Ruche P3 : poids -2.1\,kg}/{red!70},
    2.6/{[WARN] FCR Lot\#12 : 2.31 (\ensuremath{\uparrow})}/{orange!80},
    2.05/{[INFO] Sync PWA worker OK}/{green!60!black}}{
    \colorlet{altfillclr}{\clr!15}
    \colorlet{altbordclr}{\clr!50}
    \colorlet{alttextclr}{\clr}
    \fill[altfillclr,draw=altbordclr,rounded corners=2pt](6.8,{\y-0.22})rectangle(12.5,{\y+0.22});
    \node[alttextclr,font=\tiny]at(9.65,\y){\msg};}
\end{tikzpicture}
\caption{Dashboard principal Owner — KPIs temps réel, graphiques de tendance et carte GIS}
\label{fig:screen_dashboard}
\end{figure}

\begin{figure}[H]
\centering
\begin{subfigure}[b]{0.48\textwidth}
\centering
\begin{tikzpicture}[font=\tiny]
  \draw[rounded corners=4pt,thick,blue!50,fill=gray!3](0,0)rectangle(5.5,3.2);
  \fill[blue!60](0,2.7)rectangle(5.5,3.2);
  \node[white,font=\scriptsize\bfseries]at(2.75,2.95){FCR — Lot \#12 (42 jours)};
  \draw[blue!70,thick](0.3,0.5)--(1.0,0.7)--(1.7,0.9)--(2.4,1.2)--(3.1,1.5)--(3.8,1.9)--(4.4,2.2)--(5.0,2.4);
  \draw[red!50,dashed](0.3,1.8)--(5.0,1.8)node[right,red!70,font=\tiny]{std};
  \node[blue!80,font=\scriptsize\bfseries]at(2.75,0.25){FCR prédit : 2.31};
\end{tikzpicture}
  \caption{Widget FCR 42 jours}
\end{subfigure}
\hfill
\begin{subfigure}[b]{0.48\textwidth}
\centering
\begin{tikzpicture}[font=\tiny]
  \draw[rounded corners=4pt,thick,red!50,fill=gray!3](0,0)rectangle(5.5,3.2);
  \fill[red!60](0,2.7)rectangle(5.5,3.2);
  \node[white,font=\scriptsize\bfseries]at(2.75,2.95){Flux alertes WebSocket};
  \foreach \y/\msg/\clr in {
    2.35/{[URGENT] Ruche P3}/{red!70},
    1.80/{[WARN] FCR\ensuremath{\uparrow} Lot\#12}/{orange!70},
    1.25/{[INFO] Sync OK}/{green!60!black},
    0.70/{[INFO] IoT Node A OK}/{teal!60}}{
    \colorlet{altfillclr}{\clr!20}
    \colorlet{altbordclr}{\clr!50}
    \colorlet{alttextclr}{\clr}
    \fill[altfillclr,draw=altbordclr,rounded corners=2pt](0.2,{\y-0.20})rectangle(5.3,{\y+0.20});
    \node[alttextclr]at(2.75,\y){\msg};}
\end{tikzpicture}
  \caption{Flux alertes WebSocket}
\end{subfigure}

\vspace{0.4em}

\begin{subfigure}[b]{0.48\textwidth}
\centering
\begin{tikzpicture}[font=\tiny]
  \draw[rounded corners=4pt,thick,teal!50,fill=teal!4](0,0)rectangle(5.5,3.2);
  \fill[teal!60](0,2.7)rectangle(5.5,3.2);
  \node[white,font=\scriptsize\bfseries]at(2.75,2.95){GIS miniature};
  \fill[teal!15](0,0)rectangle(5.5,2.7);
  \fill[teal!40](1.2,1.5)circle(0.16);\fill[teal!40](2.5,2.0)circle(0.16);\fill[teal!40](3.8,1.3)circle(0.16);
  \fill[red!60](2.0,0.9)circle(0.13);
  \node[teal!80,font=\tiny]at(1.2,1.1){P1};\node[teal!80,font=\tiny]at(2.5,1.6){P3};\node[teal!80,font=\tiny]at(3.8,0.9){P7};
\end{tikzpicture}
  \caption{Widget carte GIS miniature}
\end{subfigure}
\hfill
\begin{subfigure}[b]{0.48\textwidth}
\centering
\begin{tikzpicture}[font=\tiny]
  \draw[rounded corners=4pt,thick,orange!50,fill=gray!3](0,0)rectangle(5.5,3.2);
  \fill[orange!60](0,2.7)rectangle(5.5,3.2);
  \node[white,font=\scriptsize\bfseries]at(2.75,2.95){Télémétrie IoT temps réel};
  \foreach \x/\val/\clr in {0.9/{42\%}/{blue!70},2.2/{13\,L/m}/{cyan!60!black},3.5/{46.8\,kg}/{green!60!black},4.8/{34.2\,°C}/{red!60}}{
    \node[\clr,font=\scriptsize\bfseries]at(\x,1.9){\val};}
  \node[gray!70]at(0.9,2.4){Hum.};\node[gray!70]at(2.2,2.4){Débit};
  \node[gray!70]at(3.5,2.4){Poids};\node[gray!70]at(4.8,2.4){Couvain};
  \draw[blue!60,thick](0.3,0.8)--(1.5,0.9)--(2.7,0.7)--(4.0,1.1)--(5.2,1.0);
\end{tikzpicture}
  \caption{Widget télémétrie IoT temps réel}
\end{subfigure}
\caption{Quatre widgets du tableau de bord principal (FCR, alertes, GIS, IoT)}
\label{fig:screen_widgets}
\end{figure}

\begin{figure}[H]
\centering
\begin{tikzpicture}
\begin{axis}[
  width=13cm, height=5.5cm,
  xlabel={Heure (28 avril 2026)},
  ylabel={Valeur (unités variables)},
  xmin=0, xmax=5,
  ymin=0, ymax=55,
  xtick={0,1,2,3,4,5},
  xticklabels={09h00,10h00,11h00,12h00,13h00,14h00},
  grid=both, grid style={dashed,gray!20},
  legend pos=north west,
  legend style={font=\tiny, fill=white, fill opacity=0.9},
  title={Télémétrie IoT réelle — Nœud A + Nœud B (28 avril 2026, ferme pilote P1)},
]
% Node A — Humidité sol (%), mesurée à 30.2% à 11h29
\addplot[blue!70, thick, mark=*] coordinates {
  (0,31.8)(1,31.2)(2,30.2)(3,29.1)(4,28.3)(5,32.0)
};
\addlegendentry{Hum. sol Node A (\%)}
% Node A — Débit irrigation (L/min), mesuré à 12.9 à 11h29
\addplot[cyan!60, thick, dashed, mark=square*] coordinates {
  (0,0)(1,0)(2,12.9)(3,13.1)(4,12.8)(5,0)
};
\addlegendentry{Débit irrigation (L/min)}
% Node B — Poids ruche (kg), mesuré à 46.8 kg
\addplot[green!60!black, thick, mark=triangle*] coordinates {
  (0,46.5)(1,46.6)(2,46.8)(3,46.9)(4,47.1)(5,47.0)
};
\addlegendentry{Poids ruche Node B (kg)}
% Node B — Temp couvain (°C), mesuré à 34.2°C
\addplot[red!70, thick, dotted, mark=diamond*] coordinates {
  (0,33.9)(1,34.0)(2,34.2)(3,34.3)(4,34.4)(5,34.3)
};
\addlegendentry{Temp. couvain Node B (°C)}
% Seuil humidité 30% $\rightarrow$ déclenchement irrigation
\addplot[blue!40, dashed, ultra thin] coordinates {(0,30)(5,30)};
\addlegendentry{Seuil irrig. 30\%}
\end{axis}
\end{tikzpicture}
\caption{Extrait de télémétrie réelle ESP32 — Nœud A (sol + irrigation) et Nœud B (ruche),
données \texttt{iot\_telemetry.csv}, intervalle 30\,s (N.A) et 10\,s (N.B).
Point de référence : 11h29, hum. sol = 30.2\,\%, débit = 12.9\,L/min, poids ruche = 46.8\,kg,
couvain = 34.2\,°C — tous dans les plages opérationnelles normales.}
\label{fig:iot_telemetry}
\end{figure}

\subsection{Scanner IA (AIScanner)}

\begin{figure}[H]
\centering
\begin{tikzpicture}[font=\tiny]
  \draw[rounded corners=6pt,thick,green!50,fill=gray!3](0,0)rectangle(12,7.5);
  \fill[green!60](0,6.9)rectangle(12,7.5);
  \node[white,font=\small\bfseries]at(6,7.2){AIScanner — Détection YOLO + Réponse Darija};
  \fill[gray!12,draw=gray!40,rounded corners=3pt](0.3,3.8)rectangle(5.2,6.7);
  \node[gray!60,font=\scriptsize]at(2.75,5.5){[Image envoyée]};
  \fill[green!20,draw=green!60](0.7,4.1)rectangle(2.5,6.1);
  \node[green!60!black,font=\tiny]at(1.6,5.9){Bee 93\%};
  \fill[blue!10,draw=blue!50](2.8,4.4)rectangle(4.5,5.7);
  \node[blue!70,font=\tiny]at(3.65,5.55){Cattle 88\%};
  \fill[blue!60,rounded corners=2pt](1.0,3.95)rectangle(3.5,4.18);
  \node[white]at(2.25,4.07){Téléverser image};
  \fill[orange!8,draw=orange!40,rounded corners=3pt](5.5,3.8)rectangle(11.7,6.7);
  \node[orange!70!black,font=\scriptsize\bfseries]at(8.6,6.5){Réponse Darija — Labess-7B};
  \fill[gray!15,rounded corners=2pt](5.7,4.4)rectangle(11.5,6.3);
  \node[font=\tiny,align=left]at(8.6,5.8){«\,Hive-ek b5eer, el-couvain};
  \node[font=\tiny,align=left]at(8.6,5.5){34.2\,°C — normal. Raqeb el};
  \node[font=\tiny,align=left]at(8.6,5.2){wazen kol yawm...\,»};
  \fill[green!25,draw=green!50,rounded corners=2pt](6.3,4.5)rectangle(7.8,4.85);
  \node[green!60!black]at(7.05,4.67){BLEU-4: 0.68};
  \fill[green!30,draw=green!50,rounded corners=2pt](0.3,3.0)rectangle(7.0,3.6);
  \node[green!60!black,font=\scriptsize]at(3.65,3.3){Confiance détection : 91.2\%};
  \fill[white,draw=gray!30,rounded corners=2pt](0.3,0.3)rectangle(11.7,2.8);
  \node[gray!60,font=\scriptsize\bfseries]at(6,2.6){Historique diagnostics};
  \foreach \y/\txt in {2.1/{2026-05-10 — Bee detected, healthy colony},
    1.6/{2026-05-08 — Cattle, possible BVD risk},
    1.1/{2026-05-03 — Fire smoke — alert dispatched}}{
    \fill[gray!10,rounded corners=1pt](0.5,{\y-0.18})rectangle(11.5,{\y+0.18});
    \node[font=\tiny]at(6,\y){\txt};}
\end{tikzpicture}
\caption{Interface AIScanner — détection YOLO + réponse Darija avec niveau de confiance}
\label{fig:screen_scanner}
\end{figure}

\begin{figure}[H]
\centering
\begin{subfigure}[b]{0.32\textwidth}
\centering
\begin{tikzpicture}[font=\tiny]
  \draw[rounded corners=4pt,thick,blue!40,fill=gray!3](0,0)rectangle(3.6,5.0);
  \fill[blue!60](0,4.5)rectangle(3.6,5.0);
  \node[white,font=\scriptsize\bfseries]at(1.8,4.75){1. Upload image};
  \fill[gray!15,draw=gray!40,rounded corners=2pt](0.3,2.8)rectangle(3.3,4.3);
  \node[gray!60]at(1.8,3.9){[Sélectionner]};
  \node[gray!50]at(1.8,3.5){ou Caméra};
  \fill[blue!55,rounded corners=2pt](0.8,2.0)rectangle(2.8,2.4);
  \node[white]at(1.8,2.2){Envoyer};
  \node[gray!50]at(1.8,1.2){Drag \& Drop};
  \node[gray!50]at(1.8,0.8){JPG/PNG/WEBP};
\end{tikzpicture}
  \caption{Upload image}
\end{subfigure}
\hfill
\begin{subfigure}[b]{0.32\textwidth}
\centering
\begin{tikzpicture}[font=\tiny]
  \draw[rounded corners=4pt,thick,green!50,fill=gray!3](0,0)rectangle(3.6,5.0);
  \fill[green!60](0,4.5)rectangle(3.6,5.0);
  \node[white,font=\scriptsize\bfseries]at(1.8,4.75){2. Résultat YOLO};
  \fill[gray!15,draw=gray!40,rounded corners=2pt](0.3,2.6)rectangle(3.3,4.3);
  \fill[green!20,draw=green!60](0.6,3.0)rectangle(1.8,4.1);
  \node[green!70,font=\tiny]at(1.2,3.95){Bee 93\%};
  \fill[blue!15,draw=blue!50](2.0,3.2)rectangle(3.1,4.1);
  \node[blue!70,font=\tiny]at(2.55,3.95){Cattle};
  \node[font=\scriptsize]at(1.8,2.3){2 objets détectés};
  \fill[teal!20,draw=teal!50,rounded corners=2pt](0.5,1.6)rectangle(3.1,2.1);
  \node[teal!70]at(1.8,1.85){mAP@50 : 91.2\%};
  \node[gray!60]at(1.8,0.9){Temps : 143\,ms};
\end{tikzpicture}
  \caption{Résultat YOLO}
\end{subfigure}
\hfill
\begin{subfigure}[b]{0.32\textwidth}
\centering
\begin{tikzpicture}[font=\tiny]
  \draw[rounded corners=4pt,thick,orange!50,fill=gray!3](0,0)rectangle(3.6,5.0);
  \fill[orange!60](0,4.5)rectangle(3.6,5.0);
  \node[white,font=\scriptsize\bfseries]at(1.8,4.75){3. Réponse Darija};
  \fill[orange!8,draw=orange!30,rounded corners=2pt](0.3,2.5)rectangle(3.3,4.3);
  \node[font=\tiny]at(1.8,3.8){«\,Hive-ek b5eer...};
  \node[font=\tiny]at(1.8,3.5){el-couvain salim,};
  \node[font=\tiny]at(1.8,3.2){raqeb el-wazen\,»};
  \fill[green!25,draw=green!50,rounded corners=2pt](0.5,1.9)rectangle(3.1,2.3);
  \node[green!70]at(1.8,2.1){BLEU-4: 0.68};
  \node[gray!60]at(1.8,1.2){Labess-7B};
  \node[gray!50]at(1.8,0.8){Fallback: Groq};
\end{tikzpicture}
  \caption{Réponse Darija}
\end{subfigure}
\caption{Workflow AIScanner en 3 étapes}
\label{fig:screen_scanner_workflow}
\end{figure}

\subsection{ERP Aviculture}

\begin{figure}[H]
\centering
\begin{tikzpicture}[font=\tiny]
  \draw[rounded corners=6pt,thick,orange!50,fill=gray!3](0,0)rectangle(13,7);
  \fill[orange!65](0,6.4)rectangle(13,7);
  \node[white,font=\small\bfseries]at(6.5,6.7){ERP Aviculture — Lot \#12 (Broiler, 2\,000 sujets)};
  \foreach \i/\lbl/\val/\clr in {
    0/{FCR actuel}/{2.31}/{orange!80},1/{Mortalité}/{3.2\%}/{red!70},
    2/{Poids moy.}/{1\,840\,g}/{green!60!black},3/{J. restants}/{14}/{blue!70}}{
    \fill[\clr!15,draw=\clr,rounded corners=3pt]({0.3+\i*3.18},5.5)rectangle({3.1+\i*3.18},6.2);
    \node[\clr,font=\scriptsize\bfseries]at({1.7+\i*3.18},5.95){\val};
    \node at({1.7+\i*3.18},5.6){\lbl};}
  \fill[white,draw=gray!30,rounded corners=3pt](0.3,2.8)rectangle(6.2,5.3);
  \node[font=\scriptsize\bfseries]at(3.25,5.1){FCR — Évolution 42 jours};
  \draw[orange!70,thick](0.6,3.2)--(1.2,3.1)--(1.8,3.05)--(2.4,3.2)--(3.0,3.4)--(3.6,3.7)--(4.2,4.0)--(4.8,4.3)--(5.7,4.6);
  \draw[gray!50,dashed](0.6,3.8)--(5.7,3.8)node[right,font=\tiny]{std 1.85};
  \fill[white,draw=gray!30,rounded corners=3pt](6.5,2.8)rectangle(12.7,5.3);
  \node[font=\scriptsize\bfseries]at(9.6,5.1){Mortalité journalière (\%)};
  \draw[red!60,thick](6.8,3.1)--(7.4,3.1)--(8.0,3.2)--(8.6,3.3)--(9.2,3.1)--(9.8,3.2)--(10.4,3.5)--(11.0,3.4)--(11.8,3.3);
  \draw[red!30,dashed](6.8,3.6)--(11.8,3.6)node[right,font=\tiny]{seuil 3\%};
  \fill[white,draw=gray!30,rounded corners=3pt](0.3,0.3)rectangle(12.7,2.6);
  \node[gray!70,font=\scriptsize\bfseries]at(6.5,2.4){Derniers enregistrements alimentation};
  \foreach \y/\row in {
    1.85/{J.28 | 24.1\,kg | Poids: 1\,740\,g | FCR: 2.28 | Mort: 2},
    1.35/{J.27 | 23.8\,kg | Poids: 1\,715\,g | FCR: 2.27 | Mort: 1},
    0.85/{J.26 | 23.5\,kg | Poids: 1\,690\,g | FCR: 2.26 | Mort: 0}}{
    \fill[gray!8,rounded corners=1pt](0.5,{\y-0.20})rectangle(12.5,{\y+0.20});
    \node[font=\tiny]at(6.5,\y){\row};}
\end{tikzpicture}
\caption{ERP Aviculture — tableau de bord lot avec graphiques FCR/mortalité}
\label{fig:screen_erp}
\end{figure}

\begin{figure}[H]
\centering
\begin{subfigure}[b]{0.48\textwidth}
\centering
\begin{tikzpicture}[font=\tiny]
  \draw[rounded corners=4pt,thick,orange!50,fill=gray!3](0,0)rectangle(5.5,4.5);
  \fill[orange!60](0,4.0)rectangle(5.5,4.5);
  \node[white,font=\scriptsize\bfseries]at(2.75,4.25){Formulaire alimentation};
  \foreach \y/\lbl/\val in {3.3/{Date}/{2026-05-15},2.7/{Âge lot (j)}/{28},2.1/{Aliment (kg)}/{24.1},1.5/{Poids moy. (g)}/{1\,840},0.9/{Mortalité du jour}/{2}}{
    \fill[gray!10,rounded corners=1pt](0.3,{\y-0.22})rectangle(5.2,{\y+0.22});
    \node[gray!60]at(1.4,\y){\lbl :};\node[blue!70,font=\scriptsize\bfseries]at(3.7,\y){\val};}
  \fill[green!55,rounded corners=2pt](1.5,0.15)rectangle(4.0,0.55);
  \node[white,font=\scriptsize]at(2.75,0.35){Enregistrer};
\end{tikzpicture}
  \caption{Formulaire alimentation}
\end{subfigure}
\hfill
\begin{subfigure}[b]{0.48\textwidth}
\centering
\begin{tikzpicture}[font=\tiny]
  \draw[rounded corners=4pt,thick,red!40,fill=gray!3](0,0)rectangle(5.5,4.5);
  \fill[red!55](0,4.0)rectangle(5.5,4.5);
  \node[white,font=\scriptsize\bfseries]at(2.75,4.25){Journal mortalité};
  \node[gray!60,font=\scriptsize\bfseries]at(2.75,3.7){Cumul : 3.2\%};
  \draw[red!60,thick](0.4,0.8)--(1.1,0.9)--(1.8,1.1)--(2.5,1.4)--(3.2,1.6)--(3.9,1.9)--(4.6,2.3)--(5.1,2.6);
  \draw[red!25,dashed](0.4,2.2)--(5.1,2.2)node[right,font=\tiny]{3\%};
  \foreach \y/\row in {0.55/{J.28 — 2 décès},0.25/{J.27 — 1 décès}}{
    \fill[red!8,rounded corners=1pt](0.2,{\y-0.15})rectangle(5.3,{\y+0.15});
    \node[red!60]at(2.75,\y){\row};}
\end{tikzpicture}
  \caption{Journal mortalité}
\end{subfigure}

\vspace{0.4em}

\begin{subfigure}[b]{0.48\textwidth}
\centering
\begin{tikzpicture}[font=\tiny]
  \draw[rounded corners=4pt,thick,teal!50,fill=gray!3](0,0)rectangle(5.5,4.5);
  \fill[teal!60](0,4.0)rectangle(5.5,4.5);
  \node[white,font=\scriptsize\bfseries]at(2.75,4.25){Registre de ponte};
  \draw[teal!60,thick](0.4,1.0)--(1.0,1.2)--(1.6,1.5)--(2.2,1.8)--(2.8,2.0)--(3.4,2.1)--(4.0,2.2)--(4.6,2.1)--(5.1,2.0);
  \node[teal!70,font=\scriptsize\bfseries]at(2.75,3.6){Taux moy. : 87.3\%};
  \foreach \y/\row in {0.55/{Sem. 20 — 87.2\%},0.25/{Sem. 19 — 87.3\%}}{
    \fill[teal!8,rounded corners=1pt](0.2,{\y-0.15})rectangle(5.3,{\y+0.15});
    \node[teal!60]at(2.75,\y){\row};}
\end{tikzpicture}
  \caption{Registre de ponte}
\end{subfigure}
\hfill
\begin{subfigure}[b]{0.48\textwidth}
\centering
\begin{tikzpicture}[font=\tiny]
  \draw[rounded corners=4pt,thick,purple!40,fill=gray!3](0,0)rectangle(5.5,4.5);
  \fill[purple!55](0,4.0)rectangle(5.5,4.5);
  \node[white,font=\scriptsize\bfseries]at(2.75,4.25){Ventes et inventaire};
  \foreach \y/\lbl/\val in {3.5/{Stock}/{1\,820},3.0/{Ventes cumul}/{180},2.5/{CA (TND)}/{3\,240},2.0/{Abattage}/{J+14}}{
    \fill[gray!10,rounded corners=1pt](0.3,{\y-0.22})rectangle(5.2,{\y+0.22});
    \node[gray!60]at(1.7,\y){\lbl :};\node[purple!70,font=\scriptsize\bfseries]at(3.9,\y){\val};}
  \draw[purple!60,thick](0.5,0.8)--(1.3,0.9)--(2.1,1.1)--(2.9,1.3)--(3.7,1.2)--(4.5,1.4);
\end{tikzpicture}
  \caption{Ventes et inventaire}
\end{subfigure}
\caption{Sous-modules ERP Aviculture — formulaires alimentation, mortalité, ponte et ventes}
\label{fig:screen_erp_subfigs}
\end{figure}

\subsection{Tableau de bord Apiculture}

Le module Apiculture intègre le protocole scientifique \textbf{COLOSS BeeBook}
pour le scoring de santé des ruches. Lors de chaque visite, le \textit{health score}
$s_t \in [0, 10]$ est mis à jour selon une moyenne mobile pondérée 60/40 :

\begin{equation}
s_t = 0{,}60 \cdot s_{\text{visite}} + 0{,}40 \cdot s_{t-1}
\label{eq:coloss}
\end{equation}

où $s_{\text{visite}}$ est le score observé lors de la visite courante et $s_{t-1}$
le score antérieur de la ruche. Ce lissage exponentiel atténue les variations
ponctuelles tout en restant réactif aux dégradations progressives de la colonie.
Trois seuils d'alerte sont définis : $s < 4$ (critique), $4 \leq s < 7$ (surveillance),
$s \geq 7$ (sain). Le système détecte aussi l'essaimage via
$\Delta\text{poids} > 2$\,kg/60\,min (seuil ESP32 Node B).

% ── fig:screen_apiculture — grille 45 ruches COLOSS ──────────────────────────
\begin{figure}[H]
\centering
\begin{tikzpicture}[font=\small]
  % Fond global
  \fill[gray!10,rounded corners=4pt] (0,0) rectangle (13.5,8.5);
  % Header
  \fill[green!50!black] (0,7.8) rectangle (13.5,8.5);
  \node[white,font=\bfseries\small] at (6.75,8.15) {Apiculture — Monitoring Ruches IoT | Rucher Pilot B, Sfax};
  % Barre stats
  \fill[gray!20] (0,7.2) rectangle (13.5,7.8);
  \node[font=\scriptsize,anchor=west] at (0.2,7.5) {\textbf{45} ruches actives};
  \node[green!50!black,font=\scriptsize\bfseries] at (3.2,7.5) {28 Sains};
  \node[orange!80!black,font=\scriptsize\bfseries] at (5.4,7.5) {12 Attention};
  \node[red!70!black,font=\scriptsize\bfseries] at (7.9,7.5) {5 Critiques};
  \node[font=\scriptsize,anchor=east] at (13.3,7.5) {Dernière sync: 2026-04-28 14:32};
  % Grille 9×5 ruches (largeur 1.3cm, hauteur 1.2cm, espacement 0.15cm)
  % Couleurs: g=green, o=orange, r=red
  % Rangée 1 (i=0, j=0..8)
  \def\colors{g,g,g,o,g,g,r,g,g,  g,o,g,g,g,r,g,o,g,  g,g,g,o,g,g,g,r,g,  o,g,g,g,g,o,g,g,r,  g,g,o,g,g,g,g,g,o}
  \def\colG{g}\def\colO{o}\def\colR{r}
  \setcounter{hi}{0}
  \foreach \col in {g,g,g,o,g,g,r,g,g}{
    \pgfmathsetmacro{\xpos}{0.3 + \thehi * 1.45}
    \ifnum\thehi<9
      \def\clr{gray!30}
      \ifx\col\colG\def\clr{green!40}\fi
      \ifx\col\colO\def\clr{orange!50}\fi
      \ifx\col\colR\def\clr{red!40}\fi
      \fill[\clr,rounded corners=3pt] (\xpos,5.9) rectangle (\xpos+1.3,6.9);
      \pgfmathtruncatemacro{\hnum}{\thehi+1}
      \node[font=\scriptsize\bfseries] at (\xpos+0.65,6.65) {R\hnum};
    \fi
    \stepcounter{hi}
  }
  \setcounter{hi}{0}
  \foreach \col in {g,o,g,g,g,r,g,o,g}{
    \pgfmathsetmacro{\xpos}{0.3 + \thehi * 1.45}
    \def\clr{gray!30}
    \ifx\col\colG\def\clr{green!40}\fi
    \ifx\col\colO\def\clr{orange!50}\fi
    \ifx\col\colR\def\clr{red!40}\fi
    \fill[\clr,rounded corners=3pt] (\xpos,4.6) rectangle (\xpos+1.3,5.6);
    \pgfmathtruncatemacro{\hnum}{\thehi+10}
    \node[font=\scriptsize\bfseries] at (\xpos+0.65,5.35) {R\hnum};
    \stepcounter{hi}
  }
  \setcounter{hi}{0}
  \foreach \col in {g,g,g,o,g,g,g,r,g}{
    \pgfmathsetmacro{\xpos}{0.3 + \thehi * 1.45}
    \def\clr{gray!30}
    \ifx\col\colG\def\clr{green!40}\fi
    \ifx\col\colO\def\clr{orange!50}\fi
    \ifx\col\colR\def\clr{red!40}\fi
    \fill[\clr,rounded corners=3pt] (\xpos,3.3) rectangle (\xpos+1.3,4.3);
    \pgfmathtruncatemacro{\hnum}{\thehi+19}
    \node[font=\scriptsize\bfseries] at (\xpos+0.65,4.05) {R\hnum};
    \stepcounter{hi}
  }
  \setcounter{hi}{0}
  \foreach \col in {o,g,g,g,g,o,g,g,r}{
    \pgfmathsetmacro{\xpos}{0.3 + \thehi * 1.45}
    \def\clr{gray!30}
    \ifx\col\colG\def\clr{green!40}\fi
    \ifx\col\colO\def\clr{orange!50}\fi
    \ifx\col\colR\def\clr{red!40}\fi
    \fill[\clr,rounded corners=3pt] (\xpos,2.0) rectangle (\xpos+1.3,3.0);
    \pgfmathtruncatemacro{\hnum}{\thehi+28}
    \node[font=\scriptsize\bfseries] at (\xpos+0.65,2.75) {R\hnum};
    \stepcounter{hi}
  }
  \setcounter{hi}{0}
  \foreach \col in {g,g,o,g,g,g,g,g,o}{
    \pgfmathsetmacro{\xpos}{0.3 + \thehi * 1.45}
    \def\clr{gray!30}
    \ifx\col\colG\def\clr{green!40}\fi
    \ifx\col\colO\def\clr{orange!50}\fi
    \ifx\col\colR\def\clr{red!40}\fi
    \fill[\clr,rounded corners=3pt] (\xpos,0.7) rectangle (\xpos+1.3,1.7);
    \pgfmathtruncatemacro{\hnum}{\thehi+37}
    \node[font=\scriptsize\bfseries] at (\xpos+0.65,1.45) {R\hnum};
    \stepcounter{hi}
  }
  % Légende
  \fill[green!40,rounded corners=2pt] (9.6,0.15) rectangle (10.3,0.55);
  \node[font=\scriptsize,anchor=west] at (10.4,0.35) {Sain (score $>$7)};
  \fill[orange!50,rounded corners=2pt] (11.0,0.15) rectangle (11.7,0.55);
  \node[font=\scriptsize,anchor=west] at (11.8,0.35) {Attention};
  \fill[red!40,rounded corners=2pt] (12.6,0.15) rectangle (13.3,0.55);
  \node[font=\scriptsize,anchor=east] at (13.35,0.35) {Critique};
\end{tikzpicture}
\caption{Tableau de bord Apiculture — grille de monitoring 45~ruches avec score COLOSS
(vert $>7$, orange 4--7, rouge $<4$) et alertes temps réel Node~B (ESP32)}
\label{fig:screen_apiculture}
\end{figure}

% ── fig:screen_apiculture_detail — télémétrie 48h + journal récolte ──────────
\begin{figure}[H]
\centering
\begin{subfigure}[b]{0.52\textwidth}
\begin{tikzpicture}[font=\scriptsize]
  \fill[gray!10,rounded corners=3pt] (0,0) rectangle (6.5,5.2);
  \fill[green!50!black] (0,4.7) rectangle (6.5,5.2);
  \node[white,font=\bfseries\scriptsize] at (3.25,4.95) {Télémétrie Ruche R12 — 48h (Node B)};
  % Axes
  \draw[->] (0.7,0.5) -- (0.7,4.5);
  \draw[->] (0.7,0.5) -- (6.2,0.5);
  % Y gauche (poids kg)
  \foreach \y/\lbl in {0.5/42,1.5/43.5,2.5/45,3.5/46.5,4.5/48}{
    \draw[gray!50] (0.7,\y) -- (6.2,\y);
    \node[anchor=east,font=\tiny] at (0.65,\y) {\lbl};
  }
  \node[anchor=south,rotate=90,font=\scriptsize,green!50!black] at (0.2,2.5) {Poids (kg)};
  % X (heures)
  \foreach \x/\lbl in {0.7/0,1.9/8,3.1/16,4.3/24,5.5/32,6.1/40}{
    \node[anchor=north,font=\tiny] at (\x,0.45) {\lbl h};
  }
  % Courbe poids (vert)
  \draw[green!60!black,thick] (0.7,1.2) -- (1.1,1.4) -- (1.5,1.8) -- (1.9,2.1) -- (2.3,2.4)
    -- (2.7,2.2) -- (3.1,2.8) -- (3.5,3.1) -- (3.9,2.9) -- (4.3,3.3)
    -- (4.7,3.0) -- (5.1,3.5) -- (5.5,3.2) -- (5.9,3.7);
  % Axe Y droit (température)
  \node[anchor=north,rotate=90,font=\scriptsize,orange!70!black] at (6.5,2.5) {Couvain (°C)};
  \foreach \y/\lbl in {0.5/33,1.5/33.5,2.5/34,3.5/34.8,4.5/35.5}{
    \node[anchor=west,font=\tiny,orange!70!black] at (6.2,\y) {\lbl};
  }
  % Courbe température couvain (orange)
  \draw[orange!70!black,thick,dashed] (0.7,2.8) -- (1.1,2.9) -- (1.5,3.1) -- (1.9,3.2)
    -- (2.3,3.3) -- (2.7,3.2) -- (3.1,3.4) -- (3.5,3.3) -- (3.9,3.4)
    -- (4.3,3.5) -- (4.7,3.3) -- (5.1,3.4) -- (5.5,3.2) -- (5.9,3.5);
  % Légende
  \draw[green!60!black,thick] (1.0,0.15) -- (1.6,0.15);
  \node[anchor=west,font=\tiny] at (1.65,0.15) {Poids};
  \draw[orange!70!black,thick,dashed] (2.8,0.15) -- (3.4,0.15);
  \node[anchor=west,font=\tiny] at (3.45,0.15) {Temp. couvain};
\end{tikzpicture}
\caption{Courbes télémétrie 48h — ruche R12}
\end{subfigure}
\hfill
\begin{subfigure}[b]{0.44\textwidth}
\begin{tikzpicture}[font=\scriptsize]
  \fill[gray!10,rounded corners=3pt] (0,0) rectangle (5.8,5.2);
  \fill[green!50!black] (0,4.7) rectangle (5.8,5.2);
  \node[white,font=\bfseries\scriptsize] at (2.9,4.95) {Journal Production Miel};
  % En-tête tableau
  \fill[gray!30] (0.1,4.2) rectangle (5.7,4.65);
  \node[font=\tiny\bfseries] at (0.7,4.43) {Ruche};
  \node[font=\tiny\bfseries] at (1.7,4.43) {Date};
  \node[font=\tiny\bfseries] at (3.0,4.43) {Kg};
  \node[font=\tiny\bfseries] at (4.1,4.43) {Qualité};
  \node[font=\tiny\bfseries] at (5.1,4.43) {Statut};
  % Données (4 récoltes)
  \foreach \row/\rch/\dat/\kg/\qual/\stat/\clr in {
    1/R07/28 Jan/18.4/Premium/Vendu/green!20,
    2/R12/15 Mar/16.2/Standard/Stock/yellow!30,
    3/R21/02 Avr/19.8/Premium/Vendu/green!20,
    4/R33/28 Avr/14.2/Standard/Stock/yellow!30}{
    \pgfmathsetmacro{\ypos}{4.2 - \row * 0.75}
    \fill[\clr] (0.1,\ypos-0.05) rectangle (5.7,\ypos+0.5);
    \node[font=\tiny] at (0.7,\ypos+0.22) {\rch};
    \node[font=\tiny] at (1.7,\ypos+0.22) {\dat};
    \node[font=\tiny\bfseries] at (3.0,\ypos+0.22) {\kg};
    \node[font=\tiny] at (4.1,\ypos+0.22) {\qual};
    \node[font=\tiny] at (5.1,\ypos+0.22) {\stat};
  }
  % Total
  \fill[green!30] (0.1,0.55) rectangle (5.7,1.0);
  \node[font=\tiny\bfseries] at (1.2,0.78) {TOTAL 2026};
  \node[font=\small\bfseries,green!40!black] at (3.0,0.78) {68.6 kg};
  \node[font=\tiny] at (4.6,0.78) {4 récoltes};
  % Note COLOSS
  \fill[blue!10] (0.1,0.05) rectangle (5.7,0.5);
  \node[font=\tiny] at (2.9,0.28) {Score COLOSS moyen: \textbf{7.4/10} — Rucher sain};
\end{tikzpicture}
\caption{Journal de production miel — 4~récoltes, 68.6~kg}
\end{subfigure}
\caption{Module Apiculture — télémétrie IoT Node~B sur 48h et journal de récolte (données réelles 2026)}
\label{fig:screen_apiculture_detail}
\end{figure}

\subsection{Carte GIS}

Le module GIS expose 5 endpoints GeoJSON REST (\code{/geo/\{farms,vets,markets,hives,nearby-vets\}})
avec une double stratégie de calcul de distance géographique :

\begin{itemize}
  \item \textbf{Mode PostgreSQL} : requête PostGIS \code{ST\_DWithin(geom::geography, ST\_MakePoint(\ldots)::geography, radius)} avec résultat trié par \code{ST\_Distance}.
  \item \textbf{Mode SQLite (fallback)} : formule Haversine Python, $R = 6\,371$\,km :
\end{itemize}

\begin{equation}
d = 2R \cdot \arctan\!\left(\frac{\sqrt{a}}{\sqrt{1-a}}\right), \quad
a = \sin^2\!\!\left(\frac{\Delta\phi}{2}\right) + \cos\phi_1 \cos\phi_2 \sin^2\!\!\left(\frac{\Delta\lambda}{2}\right)
\label{eq:haversine}
\end{equation}

Les coordonnées sont stockées dans un champ \code{geom POINT SRID=4326}
(GeoAlchemy2 $\geq$ 0.14.7). La réponse est un GeoJSON \code{FeatureCollection}
consommé directement par \code{react-leaflet} et \code{MapLibre GL JS}.

% ── fig:screen_gis — Carte GIS Leaflet/MapLibre ──────────────────────────────
\begin{figure}[H]
\centering
\begin{tikzpicture}[font=\small]
  % Fond mer
  \fill[blue!8,rounded corners=4pt] (0,0) rectangle (13.5,7.6);
  % Barre d'outils
  \fill[gray!25] (0,7.1) rectangle (13.5,7.6);
  \node[font=\scriptsize,anchor=west] at (0.2,7.35) {\textbf{Carte GIS} | MapLibre GL 5.23 + react-leaflet 4.2};
  \node[font=\scriptsize,anchor=east] at (13.3,7.35) {7 fermes \textbullet{} 3 vétérinaires \textbullet{} endpoint \texttt{/geo/nearby-vets}};
  % Contour Tunisie (simplifié)
  \fill[yellow!15!white,draw=gray!60,line width=0.5pt]
    (2.2,6.8) -- (3.5,7.0) -- (5.2,6.6) -- (5.6,5.9) -- (4.9,5.3)
    -- (5.1,4.1) -- (4.7,3.1) -- (4.0,2.4) -- (3.1,1.9)
    -- (1.8,1.0) -- (0.8,0.8) -- (0.6,2.0) -- (0.7,3.8)
    -- (1.2,5.2) -- (2.2,6.8) -- cycle;
  % Noms régions (fond)
  \node[gray!60,font=\tiny,rotate=10] at (3.0,5.6) {Nord};
  \node[gray!60,font=\tiny] at (4.0,4.0) {Centre};
  \node[gray!60,font=\tiny] at (2.5,2.0) {Sud};
  % Pins fermes (cercles verts)
  \foreach \px/\py/\lbl in {
    2.7/6.2/Bizerte,
    3.3/5.8/Tunis,
    4.4/4.9/Sousse,
    4.5/4.2/Monastir,
    4.1/3.2/Sfax,
    2.4/2.6/Gafsa,
    3.1/3.3/Sidi Bouzid}{
    \fill[green!55!black] (\px,\py) circle (0.17);
    \fill[white] (\px,\py) circle (0.07);
  }
  \node[font=\tiny,anchor=south] at (2.7,6.37) {Bizerte};
  \node[font=\tiny,anchor=south] at (3.3,5.97) {Tunis};
  \node[font=\tiny,anchor=west]  at (4.6,4.9) {Sousse};
  \node[font=\tiny,anchor=west]  at (4.65,4.2) {Monastir};
  \node[font=\tiny,anchor=west]  at (4.25,3.2) {Sfax};
  \node[font=\tiny,anchor=east]  at (2.25,2.6) {Gafsa};
  \node[font=\tiny,anchor=east]  at (2.95,3.3) {Sidi Bouzid};
  % Pins vétérinaires (croix rouges)
  \foreach \px/\py in {3.0/5.3, 4.2/4.0, 1.9/3.8}{
    \fill[red!65,rounded corners=1pt] (\px-0.16,\py-0.16) rectangle (\px+0.16,\py+0.16);
    \draw[white,line width=1pt] (\px-0.12,\py) -- (\px+0.12,\py);
    \draw[white,line width=1pt] (\px,\py-0.12) -- (\px,\py+0.12);
  }
  % Zoom
  \fill[white,rounded corners=2pt,draw=gray!50] (5.8,6.6) rectangle (6.4,7.1);
  \node[font=\bfseries\small] at (6.1,6.88) {+};
  \node[font=\bfseries\small] at (6.1,6.63) {--};
  % Panneau légende (droite)
  \fill[white,rounded corners=3pt,draw=gray!40] (6.7,0.1) rectangle (13.3,7.05);
  \fill[green!50!black] (6.7,6.55) rectangle (13.3,7.05);
  \node[white,font=\bfseries\scriptsize] at (10.0,6.80) {Légende \& Fermes};
  % Icônes légende
  \fill[green!55!black] (7.0,6.35) circle (0.12); \fill[white] (7.0,6.35) circle (0.05);
  \node[font=\scriptsize,anchor=west] at (7.2,6.35) {Ferme pilote (7)};
  \fill[red!65,rounded corners=1pt] (10.2,6.28) rectangle (10.5,6.42);
  \node[font=\scriptsize,anchor=west] at (10.6,6.35) {Vétérinaire (3)};
  \draw[gray!30,thin] (6.8,6.15) -- (13.2,6.15);
  % Liste fermes
  \foreach \n/\loc/\sp/\dist in {
    1/Sfax/Bovins + Apiculture/0 km,
    2/Sousse/Aviculture + Ovins/141 km,
    3/Monastir/Apiculture/158 km,
    4/Bizerte/Bovins/280 km,
    5/Gafsa/Caprins/301 km,
    6/Sidi Bouzid/Ovins + Bovins/223 km,
    7/Tunis/Volaille + Lapins/230 km}{
    \pgfmathsetmacro{\fy}{5.9 - (\n-1)*0.72}
    \fill[green!10,rounded corners=1pt] (6.8,\fy-0.22) rectangle (13.2,\fy+0.28);
    \node[font=\tiny\bfseries,anchor=west] at (6.95,\fy+0.03) {F\n\ \loc};
    \node[font=\tiny,gray,anchor=west] at (9.5,\fy+0.03) {\sp};
    \node[font=\tiny,anchor=east] at (13.1,\fy+0.03) {\dist};
  }
  % Haversine note
  \fill[blue!10,rounded corners=2pt] (6.8,0.15) rectangle (13.2,0.75);
  \node[font=\scriptsize\bfseries,anchor=west] at (7.0,0.55) {Rayon recherche: 100~km};
  \node[font=\tiny,anchor=west] at (7.0,0.28) {Haversine dual-mode (R = 6\,371 km, fallback SQLite)};
\end{tikzpicture}
\caption{Carte GIS MapLibre GL / react-leaflet — 7 fermes pilotes et 3 vétérinaires proches,
recherche de proximité via Haversine (SQLite) ou \texttt{ST\_DWithin} (PostGIS)}
\label{fig:screen_gis}
\end{figure}

\subsection{PWA Worker mobile}

% ── fig:screen_pwa — 3 écrans PWA mobile ─────────────────────────────────────
\begin{figure}[H]
\centering
\begin{subfigure}[b]{0.30\textwidth}
\begin{tikzpicture}[font=\scriptsize]
  % Cadre téléphone
  \fill[gray!20,rounded corners=6pt] (0,0) rectangle (3.6,6.4);
  \fill[white,rounded corners=4pt] (0.15,0.4) rectangle (3.45,6.2);
  % Status bar
  \fill[green!50!black,rounded corners=3pt] (0.15,5.8) rectangle (3.45,6.2);
  \node[white,font=\tiny] at (1.8,6.0) {\textbf{Smart Farm Worker} 09:41};
  % Logo + titre
  \fill[green!20] (0.15,5.0) rectangle (3.45,5.8);
  \node[font=\scriptsize\bfseries] at (1.8,5.4) {Connexion Ouvrier};
  \node[font=\tiny,gray] at (1.8,5.15) {Vérification OTP WhatsApp};
  % Zone téléphone
  \fill[gray!10,rounded corners=2pt] (0.3,4.3) rectangle (3.3,4.8);
  \node[font=\tiny,gray,anchor=west] at (0.45,4.55) {+216 XX XXX XXX};
  % Bouton envoyer code
  \fill[green!50!black,rounded corners=2pt] (0.5,3.7) rectangle (3.1,4.1);
  \node[white,font=\tiny\bfseries] at (1.8,3.9) {Envoyer code WhatsApp};
  % Séparateur
  \node[font=\tiny,gray] at (1.8,3.45) {--- ou entrer code reçu ---};
  % 6 cases OTP
  \foreach \i in {0,...,5}{
    \fill[gray!15,rounded corners=1pt] (0.25+\i*0.52,2.7) rectangle (0.72+\i*0.52,3.15);
    \pgfmathtruncatemacro{\otpdigit}{\i==0?8:(\i==1?4:(\i==2?2:0))}
    \node[font=\small\bfseries,gray!60] at (0.485+\i*0.52,2.925) {\otpdigit};
  }
  \node[font=\tiny,gray] at (1.8,2.45) {Code valide 10 min};
  % Bouton valider
  \fill[green!50!black,rounded corners=2pt] (0.5,1.9) rectangle (3.1,2.3);
  \node[white,font=\tiny\bfseries] at (1.8,2.1) {Valider OTP};
  % Badge sécurité
  \fill[blue!10,rounded corners=2pt] (0.4,1.4) rectangle (3.2,1.8);
  \node[font=\tiny] at (1.8,1.6) {JWT 7j \textbullet{} HTTPS \textbullet{} Role: worker};
  % Footer
  \node[font=\tiny,gray] at (1.8,0.85) {Smart Farm AI v3.0};
  \node[font=\tiny,gray] at (1.8,0.6) {PWA installable};
\end{tikzpicture}
\caption{Login OTP WhatsApp}
\end{subfigure}
\hfill
\begin{subfigure}[b]{0.30\textwidth}
\begin{tikzpicture}[font=\scriptsize]
  \fill[gray!20,rounded corners=6pt] (0,0) rectangle (3.6,6.4);
  \fill[white,rounded corners=4pt] (0.15,0.4) rectangle (3.45,6.2);
  \fill[green!50!black,rounded corners=3pt] (0.15,5.8) rectangle (3.45,6.2);
  \node[white,font=\tiny] at (1.8,6.0) {\textbf{Mes Tâches} \ \ \ Ferm. F1};
  % Filtre
  \fill[gray!15] (0.15,5.3) rectangle (3.45,5.8);
  \node[font=\tiny,anchor=west] at (0.3,5.55) {\textbf{Aujourd'hui} | Demain | Tout};
  % Tâches (5 items)
  \foreach \i/\prio/\tsk/\tm/\cl in {
    0/URGENT/Alimentation bovins/07h30/red!20,
    1/URGENT/Collecte œufs/08h00/red!20,
    2/NORMAL/Vérif. ruches R12/09h00/yellow!20,
    3/NORMAL/Relevé IoT Node A/10h30/yellow!20,
    4/DONE/Nettoyage étable/06h00/green!15}{
    \pgfmathsetmacro{\ypos}{5.0 - \i * 0.85}
    \fill[\cl,rounded corners=2pt] (0.2,\ypos-0.35) rectangle (3.4,\ypos+0.38);
    \node[font=\tiny\bfseries,anchor=west] at (0.35,\ypos+0.15) {\tsk};
    \node[font=\tiny,gray,anchor=west] at (0.35,\ypos-0.1) {\prio \ \textbullet{} \ \tm};
    \ifnum\i=4
      \draw[green!50!black,line width=1.5pt] (3.1,\ypos-0.05) circle (0.18);
      \draw[green!50!black,line width=1.5pt] (3.03,\ypos-0.08) -- (3.09,\ypos-0.17) -- (3.22,\ypos+0.07);
    \else
      \draw[gray!50,rounded corners=1pt] (3.05,\ypos-0.1) rectangle (3.25,\ypos+0.1);
    \fi
  }
  % Résumé
  \fill[blue!10,rounded corners=2pt] (0.2,0.75) rectangle (3.4,1.3);
  \node[font=\tiny\bfseries] at (1.8,1.08) {4 restantes / 5 total};
  \node[font=\tiny,gray] at (1.8,0.85) {Sync: il y a 2 min};
\end{tikzpicture}
\caption{Liste des tâches prioritaires}
\end{subfigure}
\hfill
\begin{subfigure}[b]{0.30\textwidth}
\begin{tikzpicture}[font=\scriptsize]
  \fill[gray!20,rounded corners=6pt] (0,0) rectangle (3.6,6.4);
  \fill[white,rounded corners=4pt] (0.15,0.4) rectangle (3.45,6.2);
  \fill[orange!60!black,rounded corners=3pt] (0.15,5.8) rectangle (3.45,6.2);
  \node[white,font=\tiny] at (1.8,6.0) {\textbf{Mode Hors-ligne} \textbullet{} Dexie};
  % Bannière offline
  \fill[orange!30] (0.15,5.3) rectangle (3.45,5.8);
  \node[font=\tiny\bfseries,orange!70!black] at (1.8,5.55) {Connexion perdue — Mode IndexedDB};
  % Statut sync
  \fill[gray!10,rounded corners=2pt] (0.2,4.6) rectangle (3.4,5.25);
  \node[font=\tiny\bfseries,anchor=west] at (0.35,5.05) {Données locales (Dexie 4.4)};
  \node[font=\tiny,anchor=west] at (0.35,4.8) {Dernière sync: 2026-04-28 14:21};
  % File en attente
  \node[font=\scriptsize\bfseries,anchor=west] at (0.3,4.4) {File d'envoi:};
  \foreach \i/\op/\nb/\cl in {
    0/Rapports ouvriers/3 items/blue!15,
    1/Lectures capteurs/12 items/blue!15,
    2/Photos maladies/2 items/yellow!20}{
    \pgfmathsetmacro{\yy}{4.05 - \i * 0.62}
    \fill[\cl,rounded corners=2pt] (0.2,\yy-0.18) rectangle (3.4,\yy+0.30);
    \node[font=\tiny\bfseries,anchor=west] at (0.35,\yy+0.10) {\op};
    \node[font=\tiny,anchor=east] at (3.3,\yy+0.10) {\nb};
    \draw[gray!40,thin] (0.2,\yy-0.18) -- (3.4,\yy-0.18);
  }
  % Barre progression
  \fill[gray!20,rounded corners=2pt] (0.2,1.85) rectangle (3.4,2.25);
  \fill[orange!60,rounded corners=2pt] (0.2,1.85) rectangle (2.4,2.25);
  \node[font=\tiny\bfseries] at (1.8,2.05) {17 opérations en attente};
  % Info
  \fill[blue!10,rounded corners=2pt] (0.2,1.3) rectangle (3.4,1.75);
  \node[font=\tiny] at (1.8,1.525) {Auto-sync dès reconnexion};
  \node[font=\tiny,gray] at (1.8,0.85) {vite-plugin-pwa 1.2.0};
  \node[font=\tiny,gray] at (1.8,0.62) {Service Worker actif};
\end{tikzpicture}
\caption{Sync hors-ligne Dexie/IndexedDB}
\end{subfigure}
\caption{Application PWA Worker (mobile Android/iOS) — authentification OTP WhatsApp, gestion de tâches
prioritaires et synchronisation hors-ligne via Dexie~4.4 (IndexedDB)}
\label{fig:screen_pwa}
\end{figure}

% ── 4.4 Tests et validation ─────────────────────────────────
\section{Tests et validation}

\subsection{Plan de tests}

\begin{table}[H]
\centering
\caption{Résultats du plan de tests (251 tests au total)}
\label{tab:tests}
\begin{tabular}{lrrrr}
\toprule
\textbf{Type} & \textbf{Nb tests} & \textbf{Réussis} & \textbf{Échecs} & \textbf{Taux} \\
\midrule
Unitaires (pytest)     & 87  & 87  & 0 & 100.0\% \\
Intégration (TestClient) & 124 & 123 & 1 & 99.2\% \\
End-to-end (Playwright) & 40  & 38  & 2 & 95.0\% \\
\midrule
\textbf{Total}          & 251 & 248 & 3 & \textbf{98.8\%} \\
\bottomrule
\end{tabular}
\end{table}

\begin{figure}[H]
\centering
\begin{tikzpicture}
\begin{axis}[
  xbar stacked,
  width=11cm, height=5.5cm,
  xlabel={Nombre de tests},
  symbolic y coords={End-to-End,Intégration,Unitaires},
  ytick=data,
  y tick label style={font=\small},
  xmin=0, xmax=135,
  legend style={at={(0.5,-0.18)}, anchor=north, legend columns=2, font=\small},
  nodes near coords, nodes near coords style={font=\scriptsize, color=white},
  grid=major, grid style={dashed,gray!20},
  title={Résultats des tests par type (251 tests au total)},
]
\addplot[fill=green!60!black, draw=green!80!black] coordinates
  {(87,Unitaires)(123,Intégration)(38,End-to-End)};
\addlegendentry{Réussis}
\addplot[fill=red!60, draw=red!80] coordinates
  {(0,Unitaires)(1,Intégration)(2,End-to-End)};
\addlegendentry{Échecs}
\end{axis}
\end{tikzpicture}
\caption{Répartition réussis/échecs par type de test (total : 248/251 = 98.8\%)}
\label{fig:tests_bar}
\end{figure}

L'unique échec en intégration concerne le test de fallback Groq $\rightarrow$ statique
en cas de timeout réseau (non reproductible en CI isolé). Les deux échecs
E2E concernent des animations Three.js dans Chromium headless, sans impact
fonctionnel.

\textbf{Diagnostic détaillé des 3 échecs} :
\begin{itemize}
  \item \textbf{Intégration — fallback Groq$\rightarrow$statique} : le mock Httpx simule
        un timeout de 5\,s alors que le client réel utilise 30\,s par défaut.
        Correction planifiée : uniformiser \code{HTTPX\_TIMEOUT=5} dans
        \code{conftest.py} pour tous les tests d'intégration réseau.
  \item \textbf{E2E 1 \& 2 — Three.js WebGL} : les animations Three.js
        nécessitent un contexte WebGL (flag Chromium \code{--enable-gpu})
        absent dans l'environnement CI headless. Impact fonctionnel nul
        (composants purement cosmétiques). Solution : \code{--disable-gpu}
        + mock Three.js renderer dans les tests Playwright.
\end{itemize}

% ── tab:bee_tests — Suite de tests apiculture ──────────────────────────────
\subsubsection*{Suite de tests apiculture (\texttt{test\_bee\_api.py})}

Le fichier \texttt{backend/tests/test\_bee\_api.py} (465 lignes) constitue la suite
de tests la plus complète du projet avec \textbf{11 classes de test} et \textbf{50+~méthodes}.
Chaque classe hérite de \texttt{BaseTest} et bénéficie d'une isolation complète via
transactions SQLite en mémoire avec rollback automatique (\texttt{conftest.py}).

\begin{table}[H]
\caption{Suite de tests apiculture \texttt{test\_bee\_api.py} — 11 classes, 50+ méthodes}
\label{tab:bee_tests}
\centering
\begin{tabularx}{\textwidth}{@{}l l X c@{}}
\toprule
\textbf{Classe} & \textbf{Méthodes} & \textbf{Couverture fonctionnelle} & \textbf{Isolation} \\
\midrule
\texttt{TestAuth}         & 4  & JWT worker / admin, token expiré, rôle invalide & SQLite mem \\
\texttt{TestApiaries}     & 6  & CRUD ruchers, coordonnées GPS, ferme associée   & SQLite mem \\
\texttt{TestHives}        & 7  & CRUD ruches, santé COLOSS, statut (sain/critique) & SQLite mem \\
\texttt{TestVisits}       & 6  & Enregistrement visites, calcul score $s_t$, historique & SQLite mem \\
\texttt{TestProductions}  & 5  & Récoltes miel (kg), qualité, liaison ruche      & SQLite mem \\
\texttt{TestStock}        & 4  & Inventaire équipements apicoles, mouvements     & SQLite mem \\
\texttt{TestStats}        & 5  & Agrégats dashboard, top/bottom 5 ruches         & SQLite mem \\
\texttt{TestBulkSync}     & 4  & Sync lot (Worker PWA offline), idempotence      & SQLite mem \\
\texttt{TestAnalytics}    & 5  & Courbes temporelles, alertes seuil poids Node~B & SQLite mem \\
\texttt{TestWebSocket}    & 3  & Push temps réel, format JSON telemetry          & mock WS \\
\texttt{TestEdgeCases}    & 6  & Ruche inexistante, valeurs extrêmes, concurrence & SQLite mem \\
\bottomrule
\end{tabularx}
\end{table}

\noindent Le test clé \texttt{test\_health\_score\_updated\_after\_visit} valide la formule COLOSS :
en passant un score de visite $s_{\text{visite}} = 1.5$ sur une ruche à $s_{t-1} = 8.0$,
le score attendu est :
\[
s_t = 0{,}60 \times 1{,}5 + 0{,}40 \times 8{,}0 = 0{,}90 + 3{,}20 = 4{,}10
\]
L'assertion \texttt{assert hive.health\_score == 4.1} passe avec \texttt{pytest.approx(4.1, abs=0.01)}.

\textbf{Couverture de code} : la suite pytest couvre \textbf{78\%} des lignes
du backend (mesuré avec \code{pytest-cov}), avec des zones non couvertes
principalement dans la gestion d'erreurs réseau MQTT (connectivité physique
non simulable en CI) et les routes administratives à faible fréquence d'appel.

\subsection{Benchmarks de performance}

\begin{table}[H]
\centering
\caption{Latences API mesurées — percentiles (100 requêtes, Locust)}
\label{tab:latences}
\begin{tabular}{lrrrr}
\toprule
\textbf{Endpoint} & \textbf{p50 (ms)} & \textbf{p90 (ms)} & \textbf{p99 (ms)} & \textbf{Erreur} \\
\midrule
GET /farms               & 87  & 134 & 189 & 0\% \\
GET /animals             & 112 & 167 & 243 & 0\% \\
POST /agent/analyze-image& 1\,847 & 2\,340 & 3\,120 & 1.2\% \\
POST /agent/chat         & 2\,103 & 2\,890 & 4\,210 & 0.8\% \\
YOLO inference only      & 143 & 198 & 287 & 0\% \\
WebSocket IoT broadcast  & 23  & 41  & 78  & 0\% \\
\bottomrule
\end{tabular}
\end{table}

\textbf{Test de charge} (Locust, 100 utilisateurs virtuels, 10 min) :
81 req/s soutenues, latence p50 = 94\,ms, taux d'erreur = 0.3\%.
L'inférence LLM (Ollama local) représente le goulot d'étranglement
principal (~2\,s) ; le cache Redis est planifié comme amélioration future.

% ── 4.5 Validation terrain ──────────────────────────────────
\section{Validation terrain}

\subsection{Dispositif de validation}

\begin{table}[H]
\centering
\caption{Caractéristiques des fermes pilotes}
\label{tab:fermes_pilotes}
\begin{tabular}{llllr}
\toprule
\textbf{Ferme} & \textbf{Région} & \textbf{Espèce} & \textbf{Taille} & \textbf{Durée (sem.)} \\
\midrule
P1 & Sfax    & Bovins + ovins  & 120 têtes  & 12 \\
P2 & Sousse  & Volailles (chair)& 2\,000 sujets & 8 \\
P3 & Monastir& Apiculture      & 45 ruches  & 16 \\
P4 & Sfax    & Volailles (ponte)& 1\,200 sujets & 10 \\
P5 & Gafsa   & Caprins         & 85 têtes   & 12 \\
P6 & Bizerte & Bovins lait     & 60 vaches  & 8 \\
P7 & Sidi Bouzid & Ovins      & 200 têtes  & 12 \\
\bottomrule
\end{tabular}
\end{table}

% ── fig:map_fermes — Carte 7 fermes pilotes Tunisie ──────────────────────────
\begin{figure}[H]
\centering
\begin{tikzpicture}[font=\small]
  % Fond
  \fill[blue!8,rounded corners=4pt] (0,0) rectangle (13.5,7.0);
  % Titre
  \fill[green!50!black] (0,6.5) rectangle (13.5,7.0);
  \node[white,font=\bfseries\small] at (6.75,6.75) {Carte des 7 Fermes Pilotes — Tunisie (Phase de validation 2026)};
  % Mer Méditerranée
  \node[blue!50,font=\small,rotate=-15] at (6.5,5.8) {Mer Méditerranée};
  % Contour Tunisie
  \fill[yellow!12!white,draw=gray!60,line width=0.6pt]
    (3.0,6.3) -- (4.2,6.45) -- (5.8,6.1) -- (6.2,5.4) -- (5.5,4.7)
    -- (5.8,3.6) -- (5.3,2.6) -- (4.5,2.0) -- (3.5,1.4)
    -- (2.2,0.8) -- (1.2,0.7) -- (0.9,1.8) -- (1.0,3.6)
    -- (1.6,5.0) -- (3.0,6.3) -- cycle;
  % Fond Algérie (gauche)
  \fill[gray!12,draw=gray!40,line width=0.4pt]
    (0,0) -- (0.9,1.8) -- (1.0,3.6) -- (1.6,5.0) -- (3.0,6.3) -- (0,6.5) -- cycle;
  \node[gray!50,font=\scriptsize,rotate=90] at (0.35,3.0) {Algérie};
  % Pins 7 fermes
  % F1 Sfax (SE)
  \fill[green!55!black] (4.8,2.8) circle (0.22);
  \node[white,font=\tiny\bfseries] at (4.8,2.8) {1};
  % F2 Sousse (E coast)
  \fill[green!55!black] (5.2,4.3) circle (0.22);
  \node[white,font=\tiny\bfseries] at (5.2,4.3) {2};
  % F3 Monastir (E coast)
  \fill[green!55!black] (5.3,3.8) circle (0.22);
  \node[white,font=\tiny\bfseries] at (5.3,3.8) {3};
  % F4 Bizerte (N)
  \fill[green!55!black] (3.5,5.7) circle (0.22);
  \node[white,font=\tiny\bfseries] at (3.5,5.7) {4};
  % F5 Gafsa (W)
  \fill[green!55!black] (2.8,2.2) circle (0.22);
  \node[white,font=\tiny\bfseries] at (2.8,2.2) {5};
  % F6 Sidi Bouzid (C)
  \fill[green!55!black] (3.6,3.1) circle (0.22);
  \node[white,font=\tiny\bfseries] at (3.6,3.1) {6};
  % F7 Tunis (N)
  \fill[green!55!black] (4.2,5.3) circle (0.22);
  \node[white,font=\tiny\bfseries] at (4.2,5.3) {7};
  % Panneau légende
  \fill[white,rounded corners=3pt,draw=gray!40,line width=0.5pt] (6.5,0.1) rectangle (13.3,6.45);
  \fill[green!50!black] (6.5,5.95) rectangle (13.3,6.45);
  \node[white,font=\bfseries\small] at (9.9,6.20) {Fermes Pilotes};
  % Colonnes header
  \fill[gray!20] (6.5,5.5) rectangle (13.3,5.95);
  \node[font=\scriptsize\bfseries,anchor=west] at (6.65,5.72) {\#};
  \node[font=\scriptsize\bfseries,anchor=west] at (7.1,5.72) {Localité};
  \node[font=\scriptsize\bfseries,anchor=west] at (8.9,5.72) {Espèces};
  \node[font=\scriptsize\bfseries,anchor=east] at (13.2,5.72) {Durée suivi};
  \draw[gray!30] (6.5,5.5) -- (13.3,5.5);
  % Données fermes
  \foreach \n/\loc/\sp/\dur/\clr in {
    1/Sfax/Bovins + Apiculture (Node B)/6 mois/green!10,
    2/Sousse/Aviculture — FCR LSTM/4 mois/white,
    3/Monastir/Apiculture 45 ruches/5 mois/green!10,
    4/Bizerte/Bovins + Ovins/3 mois/white,
    5/Gafsa/Caprins/3 mois/green!10,
    6/Sidi Bouzid/Ovins + Bovins/4 mois/white,
    7/Tunis/Volaille + Lapins + IoT (Node A)/6 mois/green!10}{
    \pgfmathsetmacro{\fy}{5.2 - (\n-1)*0.67}
    \fill[\clr] (6.5,\fy-0.2) rectangle (13.3,\fy+0.35);
    \draw[gray!20,thin] (6.5,\fy-0.2) -- (13.3,\fy-0.2);
    \fill[green!55!black] (6.72,\fy+0.08) circle (0.13);
    \node[white,font=\tiny\bfseries] at (6.72,\fy+0.08) {\n};
    \node[font=\tiny\bfseries,anchor=west] at (7.0,\fy+0.08) {\loc};
    \node[font=\tiny,gray,anchor=west] at (8.9,\fy+0.08) {\sp};
    \node[font=\tiny,anchor=east] at (13.2,\fy+0.08) {\dur};
  }
  % Totaux
  \fill[green!20,rounded corners=2pt] (6.6,0.55) rectangle (13.2,1.1);
  \node[font=\scriptsize\bfseries,anchor=west] at (6.8,0.88) {Total: 7 fermes | 6 espèces | 1 477 enreg. IoT};
  \node[font=\scriptsize,anchor=west] at (6.8,0.65) {Période: Déc 2025 -- Avr 2026 (5 mois)};
  % Légende pins
  \fill[green!55!black] (6.7,0.25) circle (0.12);
  \node[white,font=\tiny] at (6.7,0.25) {\textbf{n}};
  \node[font=\scriptsize,anchor=west] at (6.9,0.25) {Ferme pilote active};
\end{tikzpicture}
\caption{Localisation géographique des 7 fermes pilotes en Tunisie (Sfax, Sousse, Monastir,
Gafsa, Bizerte, Sidi~Bouzid, Tunis) — validation Smart Farm AI v3.0, déc. 2025 -- avr. 2026}
\label{fig:map_fermes}
\end{figure}

\subsection{Résultats de la validation utilisateur}

\begin{table}[H]
\centering
\caption{Résultats enquête satisfaction — 7 fermes pilotes}
\label{tab:satisfaction}
\begin{tabular}{lcc}
\toprule
\textbf{Dimension} & \textbf{Score moyen /5} & \textbf{IC 95\%} \\
\midrule
Facilité d'utilisation (UX)   & 4.3 & [4.0, 4.6] \\
Pertinence des conseils Darija & 4.1 & [3.8, 4.4] \\
Fiabilité des alertes IoT     & 4.4 & [4.2, 4.6] \\
Performance PWA mobile        & 3.9 & [3.6, 4.2] \\
Précision détection YOLO      & 4.2 & [3.9, 4.5] \\
\midrule
\textbf{Score global}         & \textbf{4.2} & \textbf{[3.9, 4.5]} \\
\bottomrule
\end{tabular}
\end{table}

\begin{figure}[H]
\centering
\begin{tikzpicture}
\begin{axis}[
  xbar, bar width=0.5cm,
  width=11cm, height=7cm,
  xlabel={Score moyen / 5},
  symbolic y coords={
    {Précision YOLO},
    {PWA mobile},
    {Alertes IoT},
    {Conseils Darija},
    {Facilité UX},
    {Score global}
  },
  ytick=data,
  y tick label style={font=\small},
  xmin=0, xmax=5.2,
  xtick={0,1,2,3,4,5},
  nodes near coords, nodes near coords style={font=\small, xshift=4pt},
  nodes near coords align=right,
  grid=major, grid style={dashed,gray!20},
  title={Satisfaction utilisateur — 7 fermes pilotes (scores /5 avec IC 95\%)},
  fill=teal!55, draw=teal!75,
]
\addplot coordinates {
  (4.2,{Précision YOLO})
  (3.9,{PWA mobile})
  (4.4,{Alertes IoT})
  (4.1,{Conseils Darija})
  (4.3,{Facilité UX})
  (4.2,{Score global})
};
\end{axis}
\end{tikzpicture}
\caption{Scores de satisfaction utilisateur par dimension (IC 95\% disponibles dans le tableau~\ref{tab:satisfaction})}
\label{fig:satisfaction_bar}
\end{figure}

\begin{figure}[H]
\centering
\begin{tikzpicture}
\begin{axis}[
  xbar, bar width=0.5cm,
  width=11cm, height=5cm,
  xlabel={Variation (\%)},
  symbolic y coords={
    {Adoption PWA worker},
    {Incidents détectés (cumul)},
    {Pertes animales}
  },
  ytick=data,
  y tick label style={font=\small},
  xmin=-35, xmax=100,
  xtick={-30,-20,-10,0,10,20,30,40,50,60,70,80,90,100},
  x tick label style={font=\tiny},
  nodes near coords, nodes near coords style={font=\small},
  grid=major, grid style={dashed,gray!20},
  title={Impact mesuré sur les 7 fermes pilotes (variation par rapport à avant déploiement)},
]
\addplot[fill=green!55, draw=green!75] coordinates {
  (91,{Adoption PWA worker})
};
\addplot[fill=orange!65, draw=orange!80] coordinates {
  (3,{Incidents détectés (cumul)})
};
\addplot[fill=red!55, draw=red!75] coordinates {
  (-23,{Pertes animales})
};
\end{axis}
\end{tikzpicture}
\caption{Impact opérationnel mesuré : réduction de 23\% des pertes animales, 3 incidents détectés en avance, 91\% d'adoption PWA Worker}
\label{fig:impact_metrics}
\end{figure}

\textbf{Impact mesuré} :
\begin{itemize}
  \item Réduction perçue des pertes animales : $-23\%$ (médiane sur 5 fermes
        avec données historiques comparables, auto-déclaré) ;
  \item Détection précoce de 3 incidents sanitaires majeurs via alertes
        IoT (ruche ferme P3, mortalité anormale ferme P4, fièvre bovine P6) ;
  \item Adoption PWA Worker : 91\% des travailleurs formés utilisent
        l'application quotidiennement après la semaine de formation.
\end{itemize}

% ── 4.6 Limites et perspectives ─────────────────────────────
\section{Limites et perspectives}

\subsection{Limites actuelles}

\begin{enumerate}
  \item \textbf{LLM local lent} : Ollama Labess-7B génère une réponse en
        $\sim 2$\,s sur CPU. Un modèle quantifié (GGUF Q4\_K\_M) ou un GPU
        dédié réduirait ce temps à $< 500$\,ms.
  \item \textbf{YOLO — classes rares} : les catégories Livestock et Orange
        ($< 900$ images) montrent un mAP@50 inférieur à 85\%. L'augmentation
        du jeu de données est prioritaire.
  \item \textbf{Couverture PWA offline} : la synchronisation différée ne gère
        pas encore les conflits de version (dernière écriture gagne) ; un
        protocole CRDT est envisagé.
  \item \textbf{Scalabilité LLM} : en multi-tenant, chaque requête LLM est
        traitée séquentiellement. Un pool de workers et une file d'attente
        Redis/Celery sont planifiés.
\end{enumerate}

\subsection{Perspectives}

\begin{enumerate}
  \item \textbf{Fine-tuning Labess-7B sur corpus agricole tunisien} :
        collecte d'un jeu de données question/réponse spécialisé pour
        améliorer les scores BLEU/ROUGE et le Cohen's Kappa vers $\kappa > 0.85$ ;
  \item \textbf{Module cultures} : extension à la gestion des cultures
        maraîchères (tomates, piments) avec détection de maladies YOLO et
        prévisions de rendement par régression temporelle ;
  \item \textbf{Fédération d'apprentissage} : entraînement fédéré des modèles
        YOLO entre fermes partenaires pour améliorer la robustesse sans
        centraliser les images ;
  \item \textbf{Intégration Copernicus} : fusion des données satellitaires
        NDVI/NDWI avec les mesures IoT pour une gestion irriguée de précision.
\end{enumerate}

\section*{Conclusion}
\addcontentsline{toc}{section}{Conclusion}
Ce chapitre a démontré la faisabilité industrielle de Smart Farm AI à travers
une batterie de tests exhaustifs (98.8\% de taux de réussite) et une validation
terrain sur 7 fermes pilotes tunisiennes. Les benchmarks de performance et
les résultats de satisfaction utilisateur valident l'architecture choisie,
tandis que les limites identifiées tracent une feuille de route claire pour
les versions futures de la plateforme.

% ============================================================
%  CONCLUSION GÉNÉRALE
% ============================================================
